% RESEARCH_DRAFT_SHA256: 1e7a741e92fb75ea81ff44662c2322fddc111005ee87e87aa0de73b248bedf7b
% Layout and source-reference wrappers expanded for the web renderer.
\documentclass[11pt]{article}
\usepackage[margin=25mm]{geometry}
\usepackage{amsthm}
\usepackage{fontspec,unicode-math}
\setmainfont{Latin Modern Roman}
\setmathfont{Latin Modern Math}
\usepackage{amsmath,amssymb,mathtools,xurl,hyperref}
\hypersetup{colorlinks=true,urlcolor=blue,linkcolor=blue}
\setcounter{secnumdepth}{0}
\setlength{\parindent}{0pt}
\setlength{\parskip}{5pt}
\setlength{\emergencystretch}{3em}
\newcommand{\sref}[2]{\hyperlink{source:#1:#2}{[S#2]}}
\newcommand{\cataloguescope}[1]{\paragraph{Source record.}#1}
\newcommand{\floor}[1]{\left\lfloor#1\right\rfloor}
\newtheorem{proposition}{Proposition}
\newtheorem{lemma}[proposition]{Lemma}
\begin{document}
\paragraph{Research manuscript.} OpenAI \texttt{gpt-6-astra}, reasoning \texttt{ultra}; 11 October 2026.
\paragraph{Preserved source body.} The complete catalogue mathematical text and sources follow with layout and source links expanded, including source qualifications and its historical preparation notes. New research begins at the section headed Research attempt.

% UnsolvedMath ID 30004014; source catalogue code OWR-16634-007
\section{Unique Games Conjecture}\label{30004014}
{\small UnsolvedMath ID 30004014 · Computer Science}
\subsection{Definitions and mathematical statement}

\paragraph{Shared notation.}Standard number systems are $\mathbb N,\mathbb Z,\mathbb Q,\mathbb R,\mathbb C$. Specialized notation and index ranges follow the cited source and the definitions below.


A Unique Games instance assigns labels from a finite alphabet to graph vertices, with each edge imposing a permutation relation on the endpoint labels. Its value is the largest fraction of satisfiable edges. The standard conjecture asks whether, for every positive tolerance, a sufficiently large alphabet makes distinguishing near-unit value from near-zero value NP-hard.
\paragraph{Record scope and qualifications.}
The catalogue title is retained as metadata. Named certificates, contracts, bounds or frameworks in that title are not asserted as results of the original source extract.
This repeats the Unique Games target already represented by canonical source records 43 and 30001374.

Source-record limitation: Duplicate Unique Games target, already present under IDs 43 and 30001374.

Related catalogue records: \href{https://mehmetmars7.github.io/Erdosproblems-llm-hunter/problem.html?type=open_problems&id=43}{43}, \href{https://mehmetmars7.github.io/Erdosproblems-llm-hunter/problem.html?type=open_problems&id=30001374}{30001374}. Separate source IDs are preserved for their provenance.
\paragraph{Statement recorded in the supplied source.}
For every $\varepsilon>0$, does there exist an alphabet size $k$ such that it is NP-hard to distinguish Unique-Games instances of value at least $1-\varepsilon$ from instances of value at most $\varepsilon$? \href{https://doi.org/10.4171/owr/2018/51}{[S2]}
\paragraph{Source status.}
The supplied catalogue labels this source record open; that metadata was checked against the supplied snapshot on 2026-10-11. The label is preserved for provenance and is not an independent verification of present open status or mathematical correctness. The catalogue literature assessment is dated 2026-08-21.
\subsection{Short English statement}
Ask whether nearly satisfiable permutation-constraint instances are hard to distinguish from ones with very low satisfiable value.
\subsection{Sources}
\begin{itemize}
\item[S1] Ulam AI and the UnsolvedMath Contributors, UnsolvedMath: A Curated Collection of Open Mathematics Problems, record 30004014 (OWR-16634-007). Catalogue curation and metadata: CC BY 4.0. The source extract is retained; definitions, English summaries and qualifications are editorial additions.
\url{https://huggingface.co/datasets/ulamai/UnsolvedMath}.
\item[S2] Complexity Theory (2019), Oberwolfach Reports, DOI 10.4171/owr/2018/51, pp. 3025--3080
\url{https://doi.org/10.4171/owr/2018/51}.
\item[S3] Theory of Computing 2025 survey discussion
\url{https://toc.cs.uchicago.edu/articles/v021a011/v021a011.pdf}.
\item[S4] A Quantum Unique Games Conjecture
\url{https://drops.dagstuhl.de/entities/document/10.4230/LIPIcs.ITCS.2025.76}.
\item Source-specific Oberwolfach Reports licences, including CC BY-SA 4.0 where applicable, are retained; see THIRD\_PARTY\_NOTICES.md.
\end{itemize}
\paragraph{Preparation and review.}
The original source and UnsolvedMath contributors retain attribution for the source extract. Codex (GPT-6; exact serving revision not exposed) assisted with definitions, English summaries, scope notes and formatting. No new human mathematical authorship or independently verified solution is asserted. Permanent source ID 30004014 and code OWR-16634-007 checked against the complete supplied registry and current repository. Distinct upstream records, including repeated questions and broad or damaged extracts, are retained for complete coverage of the source CS category. A repeated record is not a new independent conjecture. Complete mathematical-equivalence and current-literature review has not been performed. Related canonical identities: 43, 30001374. Preparation and final review follow contribution prompts C and D, with source gaps disclosed. Structural and rendering checks do not establish a new theorem or open status.
\label{end:30004014}

\section{Research attempt}
\subsection{Literature, source context and scope}
The full original catalogue body and source qualifications are preserved above. The relevant primary workshop passages were read in the cached publisher report. Historical claims and subsequent attributed results are distinguished from the mathematical work below. A restricted example, standard lemma or known theorem is not counted as new progress on the entire original target. Missing assumptions and repeated or rhetorical extracts are not silently turned into different solved problems.

\subsection{The basic pairwise marginal relaxation cannot distinguish any instance}
This is a repeated Unique Games target; its identity and provenance are retained. We test the basic local-marginal linear relaxation. For a label alphabet of size $k\ge1$, use vertex probabilities $p_v(a)$ and edge joint probabilities $q_{uv}(a,b)$, nonnegative and normalized, with each edge's marginals equal to the endpoint probabilities. Maximize the average probability that an edge permutation constraint is satisfied.
\begin{proposition}
For every nonempty-edge Unique Games instance, this relaxation has optimum 1, even when no satisfying labeling exists.
\end{proposition}
\begin{proof}
Set $p_v(a)=1/k$ for every vertex and label. For each oriented constraint $b=\pi_{uv}(a)$, put $q_{uv}(a,b)=1/k$ on its $k$ allowed pairs and zero otherwise. Because $\pi_{uv}$ is a permutation, both marginals are uniform, so the assignment is feasible. Each edge is satisfied with probability 1 in its own local table; hence the relaxation attains 1, which is also the largest possible average probability.
For a concrete integrality gap, take a triangle with two identity constraints on the binary alphabet and one flip constraint. The two identity edges would force all three labels equal, contradicting the flip, so no labeling satisfies all three. A constant labeling satisfies two edges, giving true value exactly $2/3$, while the marginal relaxation remains 1.
\end{proof}
The local tables cannot be combined into one global joint distribution supported on satisfying labelings in this example. Thus checking pairwise consistency is the exact failed inference. Stronger semidefinite or higher-order relaxations add information absent here, so this does not establish a barrier for all algorithms or settle the conjectured hardness gap.
\subsection{Verification and remaining gaps}
All eight triangle labelings and all local marginal sums were checked exactly; the values are $2/3$ and 1 as stated. The primary report discusses later progress on related label-cover gaps, but this attempt proves no new complexity reduction. It leaves the quantified near-unit versus near-zero promise unresolved and does not count a repeated record as a new independent conjecture.

\section{References and source attribution}
Primary workshop report: \url{https://doi.org/10.4171/owr/2018/51}. Relevant passages read: 2018/51 pp. 3052--3053, Unique Games conjecture, exact satisfiability qualification and related label-cover progress.

All retrievals were checked on 11 October 2026. Original question attribution remains in the preserved Sources section. External theorems are attributed and not claimed as new discoveries or fully verified merely by citation.


\paragraph{Integration review.} Codex (GPT-6; exact serving revision not exposed) reviewed the written arguments and expanded layout and source-reference wrappers. Repository metadata and the production web renderer were checked using Python 3.13 and JavaScript with MathJax 3.2.2. This model-assisted review is not independent human proof verification; no PDF compilation or formal proof checking is asserted.

\section{Completion Estimate}
\noindent\textbf{COMPLETION: 0\%}
No substantiated new progress toward the entire original target is claimed. The worked results and failed route are delimited above; a known result, illustrative model, or auxiliary obstruction is not counted as original-target completion.

\section{Final status and attribution}
\textbf{Attempt status: unresolved.} This describes this research attempt, not a new ruling on the source catalogue's status or on external literature claims.
The mathematics and this self-audit were generated by OpenAI \texttt{gpt-6-astra}, reasoning \texttt{ultra}, on 11 October 2026. Codex (GPT-6; exact serving revision not exposed) orchestrated the batch, reviewed the written mathematical arguments, and contributed formatting/repository review. This additional model review is not independent human verification. No human mathematical authorship, independent human proof checking, or formal verification is claimed. The historical preparation paragraphs in the reproduced catalogue body concern that body only.

Contribution prompt A was applied to the complete definition and supplied source record, followed by the adversarial checks of prompt B. Local reproduction files are \texttt{batch022/checks.py}, \texttt{computations.json}, and the per-ID audit JSON. They were not public links at preparation time. Numerical checks are finite and do not prove unrestricted results. No LaTeX engine was found locally; compilation was not performed by this research agent.

\end{document}
